We use the two-scale Lorenz-96 system as a small testbed for learned subgrid parametrisation and replace its 256 fast variables by three closures of the slow state: a regression polynomial, a small MLP, and the polynomial plus AR(1) red noise. All three are reimplemented, trained offline on the same data, and scored against the full two-scale truth on forecast valid time, free-run stability and the mean, variance, PDF and spatial spectrum of the slow variables, over five seeds and two scale-separation regimes ($c=10$ and $c=4$). All closures raise valid time from 0.363 (no closure) to 1.252--1.421 model time units at $c=10$ and none diverged in any run. The deterministic closures forecast best, and the registered tests do not distinguish the MLP from the polynomial when both see only the local slow variable. The AR(1) closure has the best climate at $c=10$ (spectral error 0.056 versus 0.144 for the polynomial) but a shorter valid time (1.252 versus 1.398); at $c=4$ all learned closures sit near the sampling-noise floor of the climate metrics. In three-seed ablations, giving the MLP neighbouring slow variables lowered its mean climate error (and, with seven inputs, raised valid time), the AR(1) closure no longer had the lower spectral and PDF error when deployed at forcings $F{=}18$ and $F{=}22$ (possibly because its noise amplitude was fitted at $F{=}20$; not isolated by a run), and a degree-2 polynomial gave a worse climate than a linear one despite a higher offline $R^2$. This is a small CPU study of an idealised toy system with a single 8-variable configuration and offline training only.
